\begin{notechapter}
\chapter{The Chain Rule and Automatic Differentiation}
\label{ch:chain-rule-autodiff}
\chaptersubtitle{%
  Vectors are pushed forward, covectors are pulled back}

\begin{chapterplan}
\begin{itemize}[nosep]
  \item The chain rule, with proof; the product rule for bilinear maps.
  \item Push-forward of velocities by $Df(a)$ (the Jacobian--vector product) and
        pull-back of covectors by $\dual{Df(a)}$ (the vector--Jacobian product).
  \item Forward mode and reverse mode as the two orders in which a composition
        can be evaluated; their costs.
  \item Backpropagation derived as reverse mode on a composition of layers; the
        transposes that appear are dual operators, and no inner product is
        involved.
  \item Computational graphs: fan-out corresponds to the addition of covectors.
\end{itemize}
\smallskip\noindent\textbf{Machine learning threads.}
\begin{itemize}[nosep]
  \item \texttt{torch.func.jvp}, \texttt{torch.func.vjp} and
        \texttt{torch.autograd.grad}.
  \item \texttt{loss.backward()} computes a covector; the optimizer then applies
        a Riesz map without saying so.
\end{itemize}
\end{chapterplan}
\end{notechapter}
